% SPDX-License-Identifier: CC-BY-SA-4.0
% Author: Matthieu Perrin
% Part: 
% Section: 
% Sub-section: 
% Frame: 

\begingroup

\begin{frame}{Notion de panne franche (crash)}

  \begin{alertblock}{Problèmes avec les verrous}
    \begin{itemize}
    \item Si un thread est bloqué, combien de temps va-t-il l'être ?
    \item Quid d'un thread bloqué par un thread non ordonnancé ?
    \item Vitesse d'exécution globale alignée sur le thread le plus lent
    \end{itemize}
  \end{alertblock}

  \begin{block}{Définition -- panne franche}
    Un thead \structure{en panne} n'execute plus jamais de pas de calcul
    \begin{itemize}
    \item En mémoire partagée, approximation d'un thread infiniment lent
    \item Parfois, la tolérance aux pannes est le but de la distribution
    \end{itemize}
    Un thread est dit \structure{correct} s'il ne tombe pas en panne
  \end{block}

  
  \begin{block}{Bloquant/Non-bloquant}
    \begin{itemize}
    \item Un algorithme est \structure{non-bloquant} s'il peut tolérer les pannes franches
    \item Tous les algorithmes utilisant des verrous sont bloquants
    \end{itemize}
  \end{block}

\end{frame}

\endgroup
\endinput

